% ============================================================================
% 第 1 章 绪论
% 预估页数：3-4 页
% 写作总要点：
%   - 写 "为什么做这件事"，不写技术细节
%   - 围绕 "医学问诊实训痛点 + LLM 与 Agent 时代机遇" 展开
%   - 国内外研究现状必须有真实文献引用（首次引用同时加脚注）
% ============================================================================
\chapter{绪论}

% ----------------------------------------------------------------------------
\section{研究背景与意义}
% 写作要点：
%   1. 第一段（背景）：医学教育数字化趋势 + 临床实习资源紧张 + 标准化病人
%      （SP）成本高、规范化训练机会有限
%   2. 第二段（机遇）：以 ChatGPT、GPT-4、通义千问为代表的 LLM 在医学领域
%      取得突破 \cite{singhal2023clinicalknowledge,kung2023chatgpt}
%   3. 第三段（意义）：构建一个面向医学院校的 AI 模拟诊疗系统，可解决：
%      (1) 病例数量与多样性受限；(2) 教师批量指导效率低；(3) 学情难量化
%   4. 提一句"本课题的创新点是引入多 Agent 协同与 RAG，超越简单的 LLM 套壳"

医学问诊（Medical Inquiry）连接医患沟通、临床思维与诊疗决策，是医学生临床技能训练中的基础环节。国家关于医学教育创新发展和教育信息化建设的文件均提出，要推动信息技术进入人才培养过程\cite{gov2020mededu,moe2018eduinfo}。但在实际教学中，问诊训练往往依赖课堂讲授、临床见习和标准化病人（Standardized Patient，SP）演练\cite{zhao2012sp}，资源供给并不稳定：真实病例受隐私和伦理约束，难以反复用于教学；标准化病人的招募、培训和维护成本较高；教师逐一复盘学生问诊过程也需要大量时间。随着医学生规模扩大，这些问题会直接影响课后练习次数和反馈质量。

近几年，以 GPT-4、通义千问等为代表的大语言模型（Large Language Model，LLM）在医学问答、考试评测和临床推理任务中表现出较强的语言理解能力。国内已有研究讨论 ChatGPT 对医学教育的影响\cite{qu2023chatgpt}，中文医疗大模型综述也强调，中文医疗语境下仍需专用数据集、评测体系和安全控制\cite{ji2024chinesemedllm}。HuatuoGPT\cite{zhang2023huatuogpt}、Qwen\cite{bai2023qwen}等模型的出现，使高校在中文场景中构建虚拟病人和训练反馈系统的门槛降低。国外研究中，Singhal 等\cite{singhal2023clinicalknowledge}和 Kung 等\cite{kung2023chatgpt}也从医学知识问答和考试角度验证了大模型的潜力。

不过，把通用大模型直接放进问诊训练并不可靠。单靠提示词约束时，模型有时会一边扮演患者，一边主动提示诊断或治疗方案；遇到血常规、影像描述等客观结果时，也可能生成病例中不存在的数值。问诊结束后的评价如果仍停留在自然语言建议，教师也难以形成可统计的学情数据。针对这些问题，本课题在通义千问模型之上加入多 Agent 协同、RAG、临床工具自动注入和结构化评分机制，设计并实现面向医学院校的虚拟问诊系统，使系统既能用于学生反复训练，也能为教师批改和教学分析提供数据基础。

% ----------------------------------------------------------------------------
\section{国内外研究现状}

% ----------------------------------------------------------------------------
\subsection{医学教育数字化研究现状}
% 写作要点：
%   - 国外：Harvard Medical School 的 OSCE 数字化、SP 训练系统
%   - 国内：人卫慕课、丁香医生模拟训练、各医学院校建设的虚拟仿真平台
%   - 局限：以静态脚本为主，缺乏交互能动性

医学教育数字化的探索可以追溯到 20 世纪 90 年代以计算机辅助教学为代表的电子病例库与多媒体教学系统。进入 21 世纪以来，随着 Web 技术与多媒体内容生产能力的发展，国外医学院校逐步建立了以临床综合考核（Objective Structured Clinical Examination，OSCE）数字化、虚拟病人（Virtual Patient）与高仿真模拟人（high-fidelity simulator）为代表的实训体系，部分研究还将虚拟病人与游戏化交互机制相结合，以增强学生的沉浸感和决策训练效果。在国内，随着"互联网 + 教育"政策的推进，人民卫生出版社相继建设了医学慕课和电子题库，多家医学院校依托国家级虚拟仿真实验教学项目建设了系统化的虚拟仿真实训平台，覆盖问诊、查体、急救等典型临床场景\cite{liu2021virtualsim}。

总体来看，传统数字化方案多采用预设脚本、决策树问答或知识点填空。它们的优势是流程可控、内容安全，但交互空间有限，学生很难像真实问诊那样自由追问，系统也难以根据薄弱点给出个性化反馈。因此，问诊实训需要从“内容搬到线上”进一步转向“训练过程可交互、可分析”。

% ----------------------------------------------------------------------------
\subsection{大语言模型在医学教育中的应用}
% 写作要点：
%   - Med-PaLM 在 USMLE 上达到 67\% \cite{singhal2023clinicalknowledge}
%   - ChatGPT 通过执业医师考试 \cite{kung2023chatgpt}
%   - 通义千问、ChatGLM-Med、HuatuoGPT 等中文医疗大模型
%   - 局限：直接 prompt LLM 容易"出戏"、给出诊断答案、缺乏教学层级

近三年来，大语言模型在医学教育领域的应用研究迅速增长。国内研究已经从教学模式、能力培养、风险防控等角度讨论了 ChatGPT 类生成式人工智能对医学教育的影响\cite{qu2023chatgpt}；中文医疗大模型研究则进一步围绕医疗咨询、医学影像、心理健康等任务梳理了中文医疗大模型的发展脉络与评测挑战\cite{ji2024chinesemedllm}。在模型实践层面，HuatuoGPT\cite{zhang2023huatuogpt}通过医学问答数据和医生真实数据增强模型的医疗咨询能力，Qwen\cite{bai2023qwen}系列模型则在大规模中英双语语料和指令对齐数据上进行预训练，提供了通用与医学场景兼具的开放接口。国外研究中，Singhal 等\cite{singhal2023clinicalknowledge}和 Kung 等\cite{kung2023chatgpt}也分别从临床知识编码和医学考试评测角度验证了大语言模型辅助医学教育与临床培训的潜力。

这些研究说明，大语言模型具备成为“虚拟病人”或“辅助带教工具”的基础能力。但已有工作更多停留在问答评测层面，对“问诊—检查—诊断—治疗—反馈”这一完整训练链路考虑不足。仅靠提示词也难以长期稳定地控制患者角色、客观检查结果和多轮记忆。要让系统进入常态化教学场景，还需要在模型之上增加工程化的角色分工、知识约束和评价闭环。

% ----------------------------------------------------------------------------
\subsection{AI Agent 与多智能体协作研究现状}
% 写作要点：
%   - ReAct 框架：推理 + 行动 \cite{yao2023react}
%   - Toolformer：让 LLM 学会调用工具 \cite{schick2023toolformer}
%   - LLM 自主 Agent 综述 \cite{wang2024llmagent}
%   - 多 Agent 协作（MetaGPT、AutoGen 等）在领域应用的兴起
%   - 论证：多 Agent + RAG + 工具调用 是当前 AI 应用工程的最佳实践

将大语言模型扩展为具有记忆、规划和工具使用能力的 AI Agent，是近年大模型落地应用的重要方向。Yao 等\cite{yao2023react}提出的 ReAct 框架把推理与行动交替组织起来，使模型可以在多步任务中调用外部工具；Schick 等\cite{schick2023toolformer}进一步讨论了模型学习工具调用的可能性。Lewis 等\cite{lewis2020rag}提出的 RAG 则通过外部知识检索约束生成内容，缓解幻觉和知识过时问题。这些研究为“提示词 + 检索 + 工具调用”的工程方案提供了基础。

在多智能体协作方向，AutoGen\cite{wu2023autogen}和 MetaGPT\cite{hong2024metagpt}都强调通过角色分工与对话协议完成复杂任务；Wang 等\cite{wang2024llmagent}也指出，医疗、教育等领域会出现更多专用 Agent。本课题借鉴这一思路，把标准化病人扮演、诊断评估和治疗反馈拆成不同 Agent，再由路由器按训练阶段调度，以适应医学问诊实训的教学边界。

% ----------------------------------------------------------------------------
\section{研究目标与主要工作}
% 写作要点：
%   - 研究目标：3 条简洁陈述
%       (1) 构建一套面向医学问诊实训的多角色完整 Web 系统
%       (2) 提出并实现"多 Agent + RAG + 临床工具 + 结构化评分"的医疗
%           AI 模拟诊疗框架
%       (3) 完成系统部署与功能验证
%   - 主要工作：分点列出 5--6 项

本课题的研究目标可以概括为以下三点：

第一，面向医学院校问诊实训需求，设计并实现一套覆盖学生、教师、管理员三类角色的完整 Web 化模拟诊疗系统，提供 AI 问诊训练、病例库管理、学习分析与系统管控等核心功能。

第二，在通用大语言模型之上，构建包含 Agent 编排、知识检索、会话记忆、临床工具与结构化评分的医疗 AI 模拟诊疗框架，提升问诊训练的可控性、客观性与教学反馈质量。

第三，完成系统的容器化部署与测试，检查功能流程、接口响应、安全控制和 AI Agent 效果。

围绕上述研究目标，本文的主要工作包括：

（1）需求分析与总体架构设计。梳理医学问诊实训的业务流程，提炼三类用户的核心需求，给出基于前后端分离与多 Agent 协同的整体架构、数据库 ER 模型与接口规范。

（2）多 Agent 路由的对话引擎实现。抽象 PatientAgent、DiagnosisAgent、FeedbackAgent 三类 Agent 角色，由 AgentRouter 在不同训练阶段动态选择，避免在问诊中出现角色混淆与答案泄露。

（3）轻量化 RAG、滑动窗口 Memory 与临床工具自动注入机制实现。将医学知识、对话历史与客观检查结果以工程化方式注入提示词，规避大模型在客观医学事实上的幻觉风险。

（4）多维结构化评分与规则兜底实现。以 JSON Schema 约束模型输出，给出问诊、诊断、治疗、沟通四个维度的可视化分析与改进建议，并在模型返回异常时通过规则兜底保障评分服务可用。

（5）端到端工程实现与部署测试。基于 Vue 3 与 Spring Boot 完成全栈实现，通过 Docker Compose 完成容器化部署，并从功能、性能、安全维度进行测试与验证。

% ----------------------------------------------------------------------------
\section{论文组织结构}
% 写作要点：
%   - 用一段过渡 + 编号列表说明 8 章内容
%   - 强调第 5 章为核心创新章

本论文围绕“虚拟问诊系统设计与实现”展开，全文共分为八章，按照"问题背景—相关技术—需求分析—架构设计—核心实现—模块详解—测试部署—总结展望"的逻辑层层推进。各章主要内容如下：

第 1 章 绪论。介绍课题的研究背景与意义，综述国内外在医学教育数字化、大语言模型医学应用与 AI Agent 三个方向的研究现状，并明确本文的研究目标、主要工作与组织结构。

第 2 章 相关技术概述。介绍系统涉及的关键技术，包括 Vue 3 前端框架、Spring Boot 后端框架、阿里百炼大语言模型平台、AI Agent 与多智能体协作、检索增强生成、Function Calling 工具调用、MySQL/Redis/MinIO 等数据基础设施以及 Docker 容器化部署。

第 3 章 系统需求分析。从业务、功能与非功能需求等角度分析系统需求，给出三类用户的用例图与典型用例规约。

第 4 章 系统总体设计。阐述系统的整体架构、功能模块划分、数据库 ER 模型、接口设计、RBAC 权限模型与安全架构，为后续的详细实现奠定基础。

第 5 章 AI Agent 核心技术实现。作为本论文的核心创新章节，重点呈现多 Agent 路由、RAG 知识检索、Redis 滑动窗口会话记忆、临床工具自动注入、结构化评分以及 SSE 流式响应等关键模块的设计原理与代码实现。

第 6 章 系统功能模块详细实现。详细介绍学生端、教师端、管理员端及公共模块的关键页面、交互流程与代码实现，覆盖 AI 问诊训练、病例管理、批改工作台、学习分析等业务场景。

第 7 章 系统测试与部署。给出功能、性能、安全与 AI Agent 效果验证的测试方案与结果，并描述系统基于 Docker Compose 的容器化部署过程与运行效果。

第 8 章 总结与展望。对全文工作进行总结，凝练创新点，分析现有不足并给出后续研究方向。
